% 原书第二篇第2.5节最后五题，用户上传底本印刷页83、PDF页92。
\thermocapillaryproblem{180}{83}
两颗相同的水银小液滴，
每颗直径为 $d=1.5\,\mathrm{mm}$，
在温度不变的条件下合并为一颗液滴。
求表面层自由能变化 $\Delta F$。
\exsolution
设两小液滴直径均为 $d$，
单滴球体积为
$\pi d^3/6$，
合并后体积为两倍，
故新液滴直径
$d_f=2^{1/3}d$。
单球的表面自由能为
$F=\alpha\pi d^2$。
因此初、末自由能分别为
\[
F_i=2\pi\alpha d^2,\qquad
F_f=\pi\alpha d_f^2
=\pi\alpha d^2\,2^{2/3}.
\]
自由能变化
\[
\boxed{\Delta F=\pi\alpha d^2(2^{2/3}-2)
=2\pi\alpha d^2(2^{-1/3}-1)<0}.
\]
取水银表面张力
$\alpha\simeq0.48\,\mathrm{N/m}$，
得
\[
\boxed{\Delta F\simeq-1.4\times10^{-6}\,\mathrm J
\approx-1.5\,\mathrm{\mu J}}.
\]
\excheck 合并后体积守恒而表面积减小，
所以自由能降低；
不能错误地让新液滴直径直接变为 $2d$。

\thermocapillaryproblem{181}{83}
肥皂水溶液表面张力为 $\alpha$，
外界气压为 $p_0$。
若鼓起半径为 $R$ 的肥皂泡，
求需要做的功；
注意气泡气体按可逆等温方式由外界压强
$p_0$ 的气源压入。
\exsolution
肥皂泡有两层表面，
半径从零增至 $R$ 后新增界面自由能为
\[
F_{\rm surf}=2\alpha(4\pi R^2)
=\boxed{8\pi\alpha R^2}.
\]
终态气泡内部超压为
$4\alpha/R$，
内部气体压强
$p=p_0+4\alpha/R$，
气体体积
$V=4\pi R^3/3$。
若补充考虑从外部压强 $p_0$ 的气源
把终态泡内的气体在等温方式下压至 $p$，
这批气体（其终态体积 $V$）的
可逆等温压缩功为
\[
A_{\rm gas}=pV\ln(p/p_0).
\]
因此原书所采用的全系统做功表达式为
\[
\boxed{A=8\pi\alpha R^2+
pV\ln\frac{p}{p_0}},\qquad
p=p_0+\frac{4\alpha}{R},\quad
V=\frac{4\pi R^3}{3}.
\]
\excheck 若只问膜本身增加的\emph{表面自由能}，
答案为第一项 $8\pi\alpha R^2$；
原书把从气源压缩泡内气体的功也算入，
须清楚说明所采用的功的边界定义。

\thermocapillaryproblem{182}{83}
半径为 $r$ 的肥皂泡内充有一摩尔理想气体，
外界气压为 $p_0$，
肥皂液表面张力为 $\alpha$。
求气体在泡内加热过程的摩尔热容 $C$
与定压热容 $C_p$ 的差。
\exsolution
泡内外压强差为
$4\alpha/r$，
故气体压强
$p(r)=p_0+4\alpha/r$，
体积
$V(r)=4\pi r^3/3$。
一摩尔理想气体：
\[
RT=pV=\frac{4\pi}{3}(p_0r^3+4\alpha r^2).
\]
假设表面张力在这一温区近似恒定，
对 $r$ 微分：
\[
R\frac{\mathrm dT}{\mathrm dr}
=\frac{4\pi}{3}(3p_0r^2+8\alpha r),
\qquad
\frac{\mathrm dV}{\mathrm dr}=4\pi r^2.
\]
过程热容
$C=C_V+p\,\mathrm dV/\mathrm dT$，
而 $C_p=C_V+R$，所以
\[
C-C_p
=\frac{p(4\pi r^2)}
{(4\pi/(3R))(3p_0r^2+8\alpha r)}-R.
\]
化简：
\[
\boxed{C-C_p=
\frac{4R\alpha}{3p_0r+8\alpha}
=\frac{R/2}{1+3p_0r/(8\alpha)}}.
\]
\excheck 当 $r$ 很大、毛细压差相对大气压极小时，
$C\to C_p$，符合等压过程极限。

\thermocapillaryproblem{183}{83}
把卡诺循环的思想用于一张液体薄膜，
证明在温度 $T$ 下等温地
增加单位液膜\emph{单界面面积}时，
必须交换的热量
$q=-T\,\mathrm d\alpha/\mathrm dT$，
其中 $\alpha(T)$ 是表面张力。
\exsolution
先把一层气液界面的面积记为 $S$。
液膜做可逆形变时，面积广义力为
$\alpha$，内能微分满足
$\mathrm dU=T\,\mathrm d\mathcal S+\alpha\,\mathrm dS$，
其中 $\mathcal S$ 是热力学熵。
定义自由能
$F=U-T\mathcal S$，得
\[
\mathrm dF=-\mathcal S\,\mathrm dT+\alpha\,\mathrm dS.
\]
混合二阶偏导交换：
\[
\left(\frac{\partial\mathcal S}{\partial S}\right)_T
=-\left(\frac{\partial\alpha}{\partial T}\right)_S.
\]
因此等温面积增量 $\mathrm dS>0$ 的
可逆吸热
\[
\delta Q_{\rm rev}=T\,\mathrm d\mathcal S
=-T\frac{\mathrm d\alpha}{\mathrm dT}\,\mathrm dS.
\]
单位\emph{单界面}新增面积所需热量：
\[
\boxed{q=\frac{\delta Q_{\rm rev}}{\mathrm dS}
=-T\frac{\mathrm d\alpha}{\mathrm dT}}.
\]
这一 Maxwell 推导等价于取温差
$\mathrm dT$、面积差 $\mathrm dS$ 的
无穷小可逆卡诺循环：
循环功 $\simeq-\alpha'(T)\mathrm dT\,\mathrm dS$，
由效率 $\mathrm dT/T$ 即得到同式。
\excheck 若对象为一整张肥皂薄膜，存在\emph{两个表面}，
同一投影面积增量对应的吸热应再乘2；
本题公式按单个界面定义。

\thermocapillaryproblem{184}{83}
一张肥皂薄膜在温度 $T$ 下等温地
把其投影表面积增大 $\Delta\sigma$。
肥皂水表面张力为 $\alpha(T)$，
其温度导数为 $\mathrm d\alpha/\mathrm dT$。
求表面层的：(a) 熵增 $\Delta\mathcal S$；
(b) 内能增量 $\Delta U$。
\exsolution
肥皂薄膜两侧均有气液界面，
因此总新增界面面积是 $2\Delta\sigma$。
自由能的表面增量为
$\Delta F=2\alpha\Delta\sigma$。
由
$\mathcal S=-(\partial F/\partial T)_\sigma$，
在温度固定而面积变化时，
\[
\boxed{\Delta\mathcal S
=-2\frac{\mathrm d\alpha}{\mathrm dT}
\,\Delta\sigma}.
\]
用热力学关系
$U=F+T\mathcal S$，
求同温表面层内能增量：
\[
\boxed{\Delta U
=2\left(\alpha-
T\frac{\mathrm d\alpha}{\mathrm dT}\right)
\Delta\sigma}.
\]
\excheck 面积 $\Delta\sigma$ 是薄膜的一面\emph{投影面积}，
内外两界面贡献须乘2；熵单位应为 $\mathrm{J/K}$。
